\chapter{Metodologia}
\label{cap:metodologia}

\section{Ambiente experimental}

Os experimentos foram executados em máquina com sistema operacional Linux
(\textit{kernel} 6.17), arquitetura x86-64. O código foi escrito em C++17 e
compilado com \texttt{g++} 13.3 e otimização \texttt{-O2}, sem qualquer diretiva
específica de arquitetura, de modo a manter a reprodutibilidade em outras máquinas.
A medição de tempo utiliza o relógio monotônico da máquina
(\texttt{std::chrono::steady\-\_clock}), imune a ajustes do relógio de parede
durante a execução.

\section{Organização da implementação}

As três versões exigidas pelo enunciado diferem entre si em apenas dois aspectos:
\textbf{onde a recursão é interrompida} (parâmetro $M$) e \textbf{como o pivô é
escolhido} (estratégia \texttt{PRIMEIRO}, \texttt{MEIO} ou \texttt{MEDIANA3}). Por
esse motivo, optou-se por implementar um \textbf{único motor} de ordenação,
parametrizado por esses dois valores; as três versões do enunciado são atalhos com
parâmetros fixos, conforme o Quadro \ref{tab:versoes}.

\begin{table}[H]
\centering
\caption{Parametrização das três versões a partir do motor único.}
\label{tab:versoes}
\begin{tabular}{lcc}
\hline
\textbf{Versão} & \textbf{$M$} & \textbf{Estratégia de pivô} \\
\hline
(a) Quicksort recursivo & 1 & Elemento central \\
(b) Quicksort híbrido & 40 & Elemento central \\
(c) Híbrido com mediana-de-três & 40 & Mediana de três \\
\hline
(d) Pivô inadequado (pior caso) & 1 & Primeiro elemento \\
\hline
\end{tabular}
\end{table}

Essa decisão de projeto não visa apenas à economia de código. Ela garante que as três
variantes contabilizem comparações e trocas \textbf{exatamente da mesma forma}, o que
é condição necessária para a validade da comparação experimental: se cada versão
possuísse seu próprio laço de partição, qualquer divergência na contagem seria
indistinguível de uma diferença real de desempenho. Pelo mesmo motivo, a versão (d),
usada para forçar o pior caso, é obtida apenas trocando a estratégia de pivô, sem
qualquer outra alteração no algoritmo.

\section{Convenções de contagem}

Conforme exigido pelo enunciado, os algoritmos mantêm contadores de comparações e
trocas. As convenções adotadas são as seguintes:

\begin{itemize}
  \item \textbf{Comparações}: contabiliza-se toda comparação entre \textit{chaves} do
        vetor (elemento contra pivô ou elemento contra elemento). Comparações entre
        índices, como $i \leq j$, não são contabilizadas, por constituírem controle de
        laço e não dependerem dos dados;
  \item \textbf{Trocas}: contabiliza-se toda \textit{movimentação} de chave dentro do
        vetor. No Quicksort, cada troca da partição conta como uma unidade; no
        Insertion Sort, cada deslocamento $v[j+1] \leftarrow v[j]$ também conta como
        uma unidade.
\end{itemize}

A segunda convenção merece justificativa. Contabilizar o deslocamento do Insertion
Sort como troca mantém as duas grandezas na mesma unidade semântica --- ``quantas
vezes um elemento mudou de posição'' --- permitindo que os dois algoritmos sejam
somados nas versões híbridas sem distorção. A alternativa, contar apenas permutações
de pares, subestimaria sistematicamente o trabalho realizado pelo Insertion Sort.

\section{Massas de teste}
\label{sec:massas}

Foram utilizadas cinco massas de teste, descritas na Tabela \ref{tab:massas}, com
$n$ igual a 1.000, 10.000, 100.000 e 500.000 elementos.

\begin{table}[H]
\centering
\caption{Massas de teste utilizadas.}
\label{tab:massas}
\begin{tabular}{lp{9.5cm}}
\hline
\textbf{Massa} & \textbf{Descrição} \\
\hline
Aleatório & Chaves uniformemente distribuídas em $[0, 10n)$. Representa o caso médio. \\
Ordenado & $0, 1, \ldots, n-1$. Melhor caso para pivô central; pior caso para pivô igual ao primeiro elemento. \\
Inverso & $n-1, \ldots, 1, 0$. Pior caso do Insertion Sort. \\
Repetidos & Apenas 10 chaves distintas, sorteadas uniformemente. Avalia o tratamento de chaves iguais. \\
Quase ordenado & Vetor ordenado com aproximadamente 1\% dos elementos trocados de posição. Representa a situação comum na prática. \\
\hline
\end{tabular}
\end{table}

Todas as massas são geradas de forma \textbf{determinística} a partir da tripla
(tipo, $n$, semente), utilizando o gerador Mersenne Twister com semente explícita.
Isso assegura duas propriedades: todas as versões do algoritmo são medidas sobre
exatamente os mesmos dados, e qualquer execução pode ser reproduzida integralmente.

A massa \textbf{Repetidos} merece observação. Com apenas 10 chaves distintas em até
500.000 posições, ela estressa o comportamento do algoritmo diante de chaves iguais.
O esquema de partição de Hoare adotado neste trabalho comporta-se bem nesse cenário:
quando todos os elementos de um subvetor são iguais ao pivô, os dois índices avançam
simetricamente e a partição resulta balanceada. Cabe registrar que o esquema
alternativo de Lomuto degeneraria para $\Theta(n^2)$ nessa mesma entrada --- o que
ilustra como a escolha do esquema de partição, detalhe frequentemente omitido, tem
consequências assintóticas.

\section{Procedimento de medição}
\label{sec:medicao}

O enunciado exige que cada versão seja executada várias vezes sobre as mesmas massas
de dados, a fim de obter médias confiáveis. Durante os testes preliminares,
verificou-se que essa exigência, isoladamente, é insuficiente. Duas decisões
adicionais foram necessárias e são documentadas aqui por afetarem diretamente as
conclusões.

\subsection{Descarte da fase de aquecimento}

Em medições preliminares, uma \textbf{mesma carga de trabalho} registrou
\textbf{16,2 ms} na primeira execução do processo e \textbf{5,3 ms} poucos segundos
depois --- variação de aproximadamente três vezes, causada pelo governador de
frequência do processador, e não pelo algoritmo.

O efeito não é inócuo. Como a varredura de $M$ é percorrida em ordem crescente, os
menores valores de $M$ eram medidos enquanto o processador ainda operava em
frequência reduzida. A calibração indicava, assim, $M = 100$ como valor ótimo,
enquanto o número de comparações --- métrica determinística e imune a esse efeito ---
apontava $M \approx 20$. O ``ótimo'' era artefato de aquecimento.

Cabe destacar que \textbf{repetir as execuções e tirar a média não corrige esse
viés}, por ser ele sistemático, e não aleatório: todas as repetições de uma
configuração medida cedo são igualmente penalizadas. A correção exige descartar a
fase transitória. O programa, portanto, executa uma rotina de aquecimento de dois
segundos antes de qualquer medição e descarta uma execução adicional por
configuração.

\subsection{Uso da mediana}

O ruído residual de medição --- preempção pelo escalonador do sistema operacional,
migração entre núcleos, interferência de outros processos --- só pode
\textit{acrescentar} tempo, nunca subtraí-lo. A distribuição das medidas é, portanto,
assimétrica à direita, e a média é sistematicamente puxada pelos valores atípicos.

Por essa razão, a análise apresentada no Capítulo \ref{cap:resultados} utiliza a
\textbf{mediana} dos tempos. Os arquivos de resultado registram mediana, média,
mínimo e desvio-padrão, de modo que a dispersão das medidas permanece documentada e
verificável.

\subsection{Número de repetições}

O número de repetições é adaptativo, pois entradas pequenas são ordenadas em dezenas
de microssegundos e sofrem proporcionalmente muito mais com transientes do sistema.
Nos experimentos principais utilizaram-se 10 repetições para $n \leq 10.000$, 5 para
$n = 100.000$ e 3 para $n = 500.000$. Na calibração de $M$, onde as diferenças são
mais sutis, utilizaram-se até 200 repetições para $n = 1.000$.

\subsection{Verificação de corretude}

Antes de qualquer medição, uma bateria de validação compara a saída das três versões
com a de \texttt{std::sort} em 509 casos, incluindo vetores vazios, unitários, com
todos os elementos iguais e as cinco massas em oito tamanhos distintos.
Adicionalmente, \textbf{toda} execução medida tem seu resultado verificado com
\texttt{std::is\_sorted}: uma implementação que não ordenasse corretamente
interromperia o programa, em vez de reportar um tempo de execução artificialmente
baixo.

\section{Calibração empírica de \textit{M}}
\label{sec:calibracao}

Para determinar $M$, varreu-se $M \in \{1, 2, \ldots, 200\}$ sobre as massas
Aleatório, Ordenado e Repetidos, para $n \in \{1.000;\ 10.000;\ 100.000\}$ e para as
duas estratégias de pivô.

O critério de escolha adotado \textbf{não} é o mínimo absoluto do tempo. Nas
proximidades do ótimo, a curva é praticamente plana, e o mínimo absoluto passa a ser
determinado por ruído: em execuções preliminares ele oscilou entre 12 e 100, sem
diferença real de desempenho. Adotou-se, portanto, o critério do \textbf{menor $M$
cujo tempo permanece dentro de 2\% do melhor tempo observado}, que é reprodutível
entre execuções e mantém o Insertion Sort restrito a subvetores efetivamente
pequenos.
